\section{Yoneda embeddings}\label{sec:yoneda_embeddings}

\paragraph{Yoneda embeddings}

\begin{definition}\label{def:presheaf}\mcite[def. 1.3.44]{Perrone2024StartingCategoryTheory}
  A \term{presheaf} on the \hyperref[def:category]{category} \( \cat{C} \) is a \hyperref[def:functor]{functor} from \( \cat{C}^\oppos \) to \hyperref[def:category_of_small_sets]{\( \cat{Set} \)}.
\end{definition}

\smallskip

\begin{definition}\label{def:yoneda_embeddings}\mcite[59]{Riehl2016CategoryTheory}
  \hyperref[con:currying]{Currying} the \( \hom \)-bifunctor gives us two versions of the so-called \term{Yoneda embedding}:
  \begin{paracol}{2}
    \begin{leftcolumn}
      \begin{equation*}
        \begin{aligned}
          &H_{\Anon}: \cat{C} \to [\cat{C}^\oppos, \cat{Set}], \\
          &H_A \coloneqq \hom(\anon*, A), \\
          &\underbrace{H_f}_{\mathclap{f: A \to B}} \coloneqq \seq{ \underbrace{s}_{\mathclap{C \to A}} \mapsto \underbrace{f \bincirc s}_{\mathclap{C \to B}} }_{C \in \obj(\cat{C}^\oppos)}.
        \end{aligned}
      \end{equation*}
    \end{leftcolumn}

    \begin{rightcolumn}
      \begin{equation*}
        \begin{aligned}
          &H^{\Anon}: \cat{C}^\oppos \to [\cat{C}, \cat{Set}], \\
          &H^A \coloneqq \hom(A, \anon*), \\
          &\underbrace{H^f}_{\mathclap{f: B \to A}} \coloneqq \seq{ \underbrace{s}_{\mathclap{A \to C}} \mapsto \underbrace{s \bincirc f}_{\mathclap{B \to C}} }_{C \in \obj(\cat{C})}.
        \end{aligned}
      \end{equation*}
    \end{rightcolumn}
  \end{paracol}

  Based on the variance on the functor itself (and not of the functors in its image), we call them the \term{covariant} and \term{contravariant} embedding, correspondingly.
\end{definition}
\begin{comments}
  \item Since we obtain \( H^{\Anon} \) by substituting \( \cat{C} \) with \( \cat{C}^\oppos \) in \( H_{\Anon} \), when proving properties about them, it is enough to only consider one of them. We will often find \( H_{\Anon} \) easier to work with.

  \item We show that \( H_{\Anon} \) and \( H^{\Anon} \) are \hyperref[def:functor_invertibility/injective]{injective on morphisms} (i.e. categorical embeddings) in \cref{thm:yoneda_embeddings_are_embeddings}, and show that they are \hyperref[def:functor_invertibility/full]{full functors} in \cref{thm:yoneda_embeddings_are_full}.
\end{comments}
\begin{defproof}
  We claim that \( H_f \) is a well-defined natural transformation for every morphism \( f \) in \( \cat{C} \), that \( H_{\Anon} \) is a well-defined functor.

  \SubProof{Proof of naturality of \( H_f \)} Naturality of \( H_f: \hom(\anon*, A) \to \hom(\anon*, B) \) is expressed by the following diagram, where \( p: D \to C \):
  \begin{equation*}
    \begin{aligned}
      \includegraphics[page=1]{output/def__yoneda_embedding}
    \end{aligned}
  \end{equation*}

  If \( s: C \to A \), going through the top-right corner gives us
  \begin{equation*}
    (H_f)_D(s \bincirc p) = f \bincirc (s \bincirc p),
  \end{equation*}
  while going through the bottom-left corner gives us
  \begin{equation*}
    (H_f)_C(s) \bincirc p = (f \bincirc s) \bincirc p.
  \end{equation*}

  Since composition must be associative, we conclude that the diagram commutes.

  \SubProof{Proof of functoriality of \( H_f \)}

  \SubProofOf*{def:functor/CF1} For any object \( A \) in \( \cat{C} \), \( H_{\id_A} \) is the identity natural transformation on \( \hom(A, \anon*) \).

  \SubProofOf*{def:functor/CF2} Given morphisms \( f: A \to B \) and \( g: B \to C \) in \( \cat{C} \), for every \( s: D \to A \) we have
  \begin{equation*}
    (H_{g \bincirc f})_D(s)
    =
    (g \bincirc f) \bincirc s
    =
    g \bincirc (f \bincirc s)
    =
    (H_g)_D(f \bincirc s)
    =
    (H_g)_D( (H_f)_D(s) ),
  \end{equation*}
  thus
  \begin{equation*}
    H_f(g \bincirc f) = H_g \bincirc H_f.
  \end{equation*}
\end{defproof}

\begin{proposition}\label{thm:yoneda_embeddings_are_embeddings}
  The \hyperref[def:yoneda_embeddings]{Yoneda embeddings} are \hyperref[def:functor_invertibility/injective]{categorical embedding} (i.e. is injective on morphisms).
\end{proposition}
\begin{proof}
  It will be sufficient to consider the covariant embedding \( H_{\Anon*} \). Suppose that \( H_f = H_{f'} \) for two morphisms \( f: A \to B \) and \( f': A' \to B' \) in \( \cat{C} \).

  Then \( H_A = \hom(\anon*, A) = \hom(\anon*, A') = H_{A'} \), which is only possible if \( A = A' \) since \( \hom(C, A) \) and \( \hom(C, A') \) would otherwise be disjoint, and similarly \( B = B' \).

  Thus, we have parallel morphisms \( f, f': A \to B \), and, for every \( s: C \to A \), we have
  \begin{equation*}
    (H_f)_C(s)
    =
    f \bincirc s
    =
    f' \bincirc s
    =
    (H_f')_C(s).
  \end{equation*}

  In particular, for \( s = \id_A \), we have \( f = g \).
\end{proof}

\begin{remark}\label{rem:yoneda_embedding_variance}
  Both the covariant and contravariant Yoneda embeddings from \cref{def:yoneda_embeddings} are presented in the literature, although generally \enquote{the Yoneda embedding} may refer to either.

  The name \enquote{covariant (resp. contravariant) Yoneda embedding} is also used in \cite[thm. 2.2.4]{Riehl2016CategoryTheory}, while the exact symbols \( H_{\Anon} \) and \( H^{\Anon} \) are used in \cite[def. 4.1.15; def. 4.1.21]{Leinster2014BasicCategoryTheory}. The symbol \( h_A \) (instead of our \( H^A \)) is used for the covariant embedding in \cite[695]{Eisenbud1995CommutativeAlgebra}, while \( y \) is used in \cite[def. 8.1]{Awodey2010CategoryTheory}. The symbols \( Y \) and \( Y' \) are also used for the covariant and contravariant embedding, respectively, in \cite[61]{MacLane1998CategoryTheory}.

  Additionally, this leads to two variants of \fullref{thm:yonedas_lemma}. We present both, but generally different authors prefer different versions. For example, the lemma for covariant functors that uses the contravariant embedding \( H^{\Anon} \) is presented in
  \cite[61]{MacLane1998CategoryTheory},
  \cite[lemma 4.2.1]{Schubert1972Categories} and
  \cite[thm. 2.2.4]{Riehl2016CategoryTheory},
  \cite[lemma A5.1]{Eisenbud1995CommutativeAlgebra},
  while the lemma for contravariant functors that uses \( H_{\Anon} \) is presented in
  \cite[lemma 8.2]{Awodey2010CategoryTheory},
  \cite[lemma 2.2.1]{Perrone2024StartingCategoryTheory} and
  \cite[thm. 4.2.1]{Leinster2014BasicCategoryTheory}.
\end{remark}

\begin{proposition}\label{thm:isomorphisms_via_yoneda_embedding}
  In a \hyperref[def:small_category]{locally small} \hyperref[def:category]{category} \( \cat{C} \), the objects \( A \) and \( B \) are \hyperref[def:morphism_invertibility/isomorphism]{isomorphic} if and only if the functors \( \hom(A, \anon*) \) and \( \hom(B, \anon*) \) (or \( \hom(\anon*, A) \) and \( \hom(\anon*, B) \)) are \hyperref[def:natural_isomorphism]{naturally isomorphic}.
\end{proposition}
\begin{proof}
  We will focus on the \hyperref[def:yoneda_embeddings]{covariant Yoneda embedding} \( H_{\Anon*}: \cat{C} \to [\cat{C}^\oppos, \cat{Set}] \), corresponding to \( \hom(\anon*, A) \) and \( \hom(\anon*, B) \).

  \SufficiencySubProof Fix an isomorphism \( f: A \to B \) in \( \cat{C} \). The Yoneda embedding \( {H_{\Anon*}: \cat{C} \to [\cat{C}^\oppos, \cat{Set}]} \) gives us natural transformations \( H_f: \hom(\anon*, A) \Rightarrow \hom(\anon*, B) \), which is an isomorphism by \cref{thm:def:functor/two_sided_inverses}.

  \NecessitySubProof Fix a natural isomorphism \( \varphi: \hom(\anon*, A) \Rightarrow \hom(\anon*, B) \).

  Denote \( \varphi_A(\id_A): A \to B \) by \( f \) and \( \varphi_B^{-1}(\id_B): B \to A \) by \( g \). We claim that \( g \) is the inverse of \( f \).

  Indeed, the naturality of \( \varphi \) implies that the following diagram commutes:
  \begin{equation*}
    \begin{aligned}
      \includegraphics[page=1]{output/thm__isomorphisms_via_yoneda_embedding}
    \end{aligned}
  \end{equation*}

  Then
  \begin{equation*}
    \underbrace{\varphi_B(\id_A \bincirc g)}_{\varphi_B(\varphi_B^{-1}(\id_B))} = \underbrace{\varphi_A(\id_A)}_f \bincirc g,
  \end{equation*}
  thus \( g \) is a right inverse of \( f \).

  Similarly, the following diagram commutes:
  \begin{equation*}
    \begin{aligned}
      \includegraphics[page=2]{output/thm__isomorphisms_via_yoneda_embedding}
    \end{aligned}
  \end{equation*}

  Then
  \begin{equation*}
    \underbrace{\varphi_A^{-1}(\id_B \bincirc f)}_{\varphi_A^{-1}(\varphi_A(\id_A))} = \underbrace{\varphi_B^{-1}(\id_B)}_g \bincirc f,
  \end{equation*}
  thus \( g \) is a left inverse of \( f \).

  Therefore, \( g \) is the two-sided inverse of \( f \).
\end{proof}

\paragraph{Yoneda's lemma}

\begin{lemma}[Yoneda's lemma]\label{thm:yonedas_lemma}
  Fix a \hyperref[def:small_category]{locally small} \hyperref[def:category]{category} \( \cat{C} \). We will state two related results, one for \hyperref[def:functor]{covariant} and one for \hyperref[def:contravariant_functor]{contravariant functors} from \( \cat{C} \) into \( \cat{Set} \).

  \begin{thmenum}
    \thmitem{thm:yonedas_lemma/component}
    \begin{paracol}{2}
      \begin{leftcolumn}
        For every object \( A \) and every functor \( {F: \cat{C} \to \cat{Set}} \), every value \( x \) from \( F(A) \) gives rise to the \hyperref[def:natural_transformation]{natural transformation}
        \begin{equation*}
          \begin{aligned}
            &\alpha_{(A, F)}(x): \hom(A, \anon*) \Rightarrow F, \\
            &\alpha_{(A, F)}(x) \coloneqq \seq[\big]{ \underbrace{s}_{\mathclap{A \to C}} \mapsto \underbrace{F(s)}_{\mathclap{\qquad F(A) \to F(C)}}(x) }_{C \in \obj(\cat{C})}.
          \end{aligned}
        \end{equation*}
      \end{leftcolumn}

      \begin{rightcolumn}
        For every object \( A \) and every presheaf \( {F: \cat{C}^{\oppos} \to \cat{Set}} \), every value \( y \) from \( F(A) \) gives rise to the natural transformation
        \begin{equation*}
          \begin{aligned}
            &\alpha'_{(A, F)}(y): \hom(\anon*, A) \Rightarrow F, \\
            &\alpha'_{(A, F)}(y) \coloneqq \seq[\big]{ \underbrace{s}_{\mathclap{C \to A}} \mapsto \underbrace{F(s)}_{\mathclap{\qquad F(A) \to F(C)}}(y) }_{C \in \obj(\cat{C}^\oppos)}.
          \end{aligned}
        \end{equation*}
      \end{rightcolumn}
    \end{paracol}

    \thmitem{thm:yonedas_lemma/transformation}
    \begin{paracol}{2}
      \begin{leftcolumn}
        Moreover, \( \alpha_{(A, F)} \) extends to a \hyperref[def:natural_transformation]{natural transformation} \( \alpha: L \to R \) between the functors
        \begin{align*}
          &L: \cat{C} \times [\cat{C}, \cat{Set}] \to \cat{Set}, \\
          &L(A, F) \coloneqq F(A), \\
          &L(\underbrace{f}_{A \to B}, \underbrace{\varphi}_{F \Rightarrow G}) \coloneqq \underbrace{\varphi_B \bincirc F(f)}_{\T*{or} G(f) \bincirc \varphi_A}
        \intertext{and}
          &R: \cat{C} \times [\cat{C}, \cat{Set}] \to \cat{Set}, \\
          &R(A, F) \coloneqq \hom(H^A, F), \\
          &R(\underbrace{f}_{A \to B}, \underbrace{\varphi}_{F \Rightarrow G}) \coloneqq \parens[\big]{ \underbrace{\psi}_{\mathclap{H^A \Rightarrow F}} \mapsto \underbrace{\varphi \bincirc \psi \bincirc H^f}_{H^B \Rightarrow G}}.
        \end{align*}
      \end{leftcolumn}

      \begin{rightcolumn}
        Moreover, \( \alpha'_{(A, F)} \) extends to a \hyperref[def:natural_isomorphism]{natural isomorphism} \( \alpha: L \to R \) between the functors
        \begin{align*}
          &L': \cat{C}^\oppos \times [\cat{C}^\oppos, \cat{Set}] \to \cat{Set}, \\
          &L'(A, F) \coloneqq F(A), \\
          &L'(\underbrace{f}_{B \to A}, \underbrace{\varphi}_{F \Rightarrow G}) \coloneqq \underbrace{\varphi_B \bincirc F(f)}_{\T*{or} G(f) \bincirc \varphi_A}
        \intertext{and}
          &R': \cat{C}^\oppos \times [\cat{C}^\oppos, \cat{Set}] \to \cat{Set}, \\
          &R'(A, F) \coloneqq \hom(H_A, F), \\
          &R'(\underbrace{f}_{B \to A}, \underbrace{\varphi}_{F \Rightarrow G}) \coloneqq \parens[\big]{ \underbrace{\psi}_{\mathclap{H_A \Rightarrow F}} \mapsto \underbrace{\varphi \bincirc \psi \bincirc H_{f}}_{H_B \Rightarrow G}}.
        \end{align*}
      \end{rightcolumn}
    \end{paracol}

    \thmitem{thm:yonedas_lemma/isomorphism}
    \begin{paracol}{2}
      \begin{leftcolumn}
        The transformation \( \alpha: L \to R \) is a \hyperref[def:natural_isomorphism]{natural isomorphism} whose inverse \( {\beta: R \to L} \) has components
        \begin{equation*}
          \beta_{(A, F)}(\psi: H^A \Rightarrow F) \coloneqq \underbrace{\psi_A(\id_A)}_{F(A)}.
        \end{equation*}
      \end{leftcolumn}

      \begin{rightcolumn}
        The transformation \( \alpha': L' \to R' \) is a \hyperref[def:natural_isomorphism]{natural isomorphism} whose inverse \( {\beta': R' \to L'} \) has components
        \begin{equation*}
          \beta'_{(A, F)}(\psi: H_A \Rightarrow F) \coloneqq \underbrace{\psi_A(\id_A)}_{F(A)}.
        \end{equation*}
      \end{rightcolumn}
    \end{paracol}
  \end{thmenum}
\end{lemma}
\begin{comments}
  \item We obtain \( L' \) and \( R' \) by replacing \( \cat{C} \) with \( \cat{C}^\oppos \) in \( L \) and \( R \).
\end{comments}
\begin{proof}
  We will prove the lemma for covariant functors, i.e. we will show that \( \alpha: L \Rightarrow R \) is a natural isomorphism. The analogous statement for \( \alpha': L' \Rightarrow R' \) follows directly.

  \SubProofOf{thm:yonedas_lemma/component} Denote \( \alpha_{(A, F)}(x) \) by \( \psi \). Given \( g: C \to D \) in \( \cat{C} \), we must show that the following diagram commutes:
  \begin{equation*}
    \begin{aligned}
      \includegraphics[page=1]{output/thm__yonedas_lemma}
    \end{aligned}
  \end{equation*}

  For a fixed morphism \( s: A \to C \), going through the top-right corner gives us
  \begin{equation*}
    \psi_D(g \bincirc s)
    =
    F(g \bincirc s)(x)
  \end{equation*}
  while going through the bottom-left corner gives us
  \begin{equation*}
    F(g)( \psi_C(s) )
    =
    F(g)( F(s)(x) )
    =
    [F(g) \bincirc F(s)](x)
    \reloset {\ref{def:functor/CF2}} =
    F(g \bincirc s)(x).
  \end{equation*}

  \SubProofOf{thm:yonedas_lemma/transformation} Given objects \( A \) and \( B \) in \( \cat{C} \), presheaves \( F, G: \cat{C}^\oppos \to \cat{Set} \) and a natural transformation \( \varphi: F \to G \), we must show that the following diagram commutes:
  \begin{equation*}
    \begin{aligned}
      \includegraphics[page=2]{output/thm__yonedas_lemma}
    \end{aligned}
  \end{equation*}

  For any element \( x \) of \( F(A) \), going through the top-right corner gives us
  \begin{equation*}
    \alpha_{(B, G)}\parens[\big]{ L(f, \varphi)(x) }
    =
    \alpha_{(B, G)}\parens[\big]{ \varphi_B\parens[\big]{ F(f)(x) } }
    =
    \seq[\big]{ \underbrace{t}_{\mathclap{B \to C}} \mapsto G(t)\parens[\big]{ \varphi_B\parens[\big]{ F(f)(x) } } }_{C \in \obj(\cat{C})}.
  \end{equation*}

  Because of the naturality of \( \varphi: F \to G \), for a fixed morphism \( t: B \to C \), it follows that
  \begin{equation*}
    G(t) \bincirc \varphi_B
    =
    \varphi_C \bincirc F(t),
  \end{equation*}
  thus,
  \begin{equation*}
    \alpha_{(B, G)}\parens[\big]{ L(f, \varphi)(x) }_C(t)
    =
    \varphi_C\parens[\big]{ \parens[\big]{ F(t) \bincirc F(f) }(x) }
    =
    \varphi_C\parens[\big]{ F(t \bincirc f)(x) }.
  \end{equation*}

  Going through the bottom-left corner gives us the same value:
  \begin{align*}
    R(f, \varphi)\parens[\big]{ \alpha_{(A, F)}(x) }
    &=
    R(f, \varphi)\parens[\big]{ \seq{ \underbrace{s}_{\mathclap{A \to C}} \mapsto F(s)(x) }_{C \in \obj(\cat{C})} }
    = \\ &=
    \seq[\big]{ \varphi_C \bincirc \parens[\big]{ \underbrace{s}_{\mathclap{A \to C}} \mapsto F(s)(x) } \bincirc \underbrace{(H^f)_C}_{\mathclap{t \mapsto t \bincirc f}} }_{C \in \obj(\cat{C})}
    = \\ &=
    \seq[\big]{ \varphi_C \bincirc \parens[\big]{ \underbrace{t}_{\mathclap{B \to C}} \mapsto F(t \bincirc f)(x) } }_{C \in \obj(\cat{C})}
    = \\ &=
    \seq[\big]{ \underbrace{t}_{\mathclap{B \to C}} \mapsto \varphi_C\parens[\big]{ F(t \bincirc f)(x) } }_{C \in \obj(\cat{C})}.
  \end{align*}

  \SubProofOf{thm:yonedas_lemma/isomorphism} It is straightforward to verify that \( \beta_{(A, F)} \) is a left inverse of \( \alpha_{(A, F)} \):
  \begin{equation*}
    \beta_{(A, F)}(\alpha_{(A, F)}(x))
    =
    F(\id_A)(x)
    \reloset {\ref{def:functor/CF1}} =
    \id_{F(A)}(x)
    =
    x
  \end{equation*}

  The right inverse is more subtle to verify. We have
  \begin{equation*}
    \alpha_{(A, F)}(\beta_{(A, F)}(\psi))
    =
    \alpha_{(A, F)}(\psi_A(\id_A))
    =
    \seq[\big]{ s \mapsto F(s)(\psi_A(\id_A)) }_{C \in \obj(\cat{C})}
  \end{equation*}

  The naturality of \( \psi \) implies that, for any \( s: A \to C \),
  \begin{equation*}
    F(s) \bincirc \psi_A = \psi_C \bincirc H^A(s),
  \end{equation*}
  thus
  \begin{equation*}
    F(s)\parens[\big]{ \psi_A(\id_A) } = \psi_C\parens[\big]{ H^A(s)(\id_A) }.
  \end{equation*}

  But \( H^A(s) \) is a function that maps \( t: A \to A \) to \( s \bincirc t \), thus
  \begin{equation*}
    H^A(s)(\id_A) = s
  \end{equation*}
  and
  \begin{equation*}
    \alpha_{(A, F)}(\beta_{(A, F)}(\psi))
    =
    \seq[\big]{ s \mapsto \varphi_C(s) }_{C \in \obj(\cat{C})}
    =
    \varphi.
  \end{equation*}

  We have thus shown that \( \beta_{(A, F)} \) is both a left and right inverse of \( \alpha_{(A, F)} \). By \cref{thm:natural_isomorphism}, \( \beta: R \to L \) is a natural transformation and a two-sided inverse of \( \alpha: L \to R \).
\end{proof}

\begin{theorem}[Yoneda's embedding theorem]\label{thm:yonedas_embedding_theorem}\mcite[thm. 2.1.16]{Perrone2024StartingCategoryTheory}
  In a \hyperref[def:small_category]{locally small} \hyperref[def:category]{category}, the \( \hom \)-bifunctor \( \hom(A, B) \) is \hyperref[def:natural_isomorphism]{naturally isomorphic} to both
  \begin{equation*}
    \hom(\hom(A, \anon*), \hom(B, \anon*))
  \end{equation*}
  and
  \begin{equation*}
    \hom(\hom(\anon*, A), \hom(\anon*, B)).
  \end{equation*}
\end{theorem}
\begin{proof}
  For every object \( A \) in \( \cat{C} \) and every functor \( F: \cat{C} \to \cat{Set} \), \fullref{thm:yonedas_lemma} gives us a natural isomorphism \( \alpha: L \to R \) between \( L(A, F) = F(A) \) and \( {R(A, F) = \hom(H^A, F)} \).

  Consider the auxiliary functor
  \begin{equation*}
    \begin{aligned}
      &P: \cat{C} \times \cat{C} \to \cat{C} \times [\cat{C}, \cat{Set}], \\
      &P(A, B) \coloneqq (A, H^B), \\
      &P(f: A \to A', g: B \to B') \coloneqq (f, H^{g^\oppos}). \\
    \end{aligned}
  \end{equation*}

  Because of the naturality of \( \alpha \), the family
  \begin{equation*}
    \seq{ \alpha_{(A, H^B)} }_{(A, B) \in \obj(\cat{C})^2}
  \end{equation*}
  is a natural transformation from \( L \bincirc P \) and \( R \bincirc P \).

  We can similarly handle the case for the contravariant \( \hom \)-functors.
\end{proof}

\begin{corollary}\label{thm:yoneda_embeddings_are_full}
  The \hyperref[def:yoneda_embeddings]{Yoneda embeddings} are \hyperref[def:functor_invertibility/full]{full functors}.
\end{corollary}
\begin{proof}
  Follows from \cref{thm:yonedas_embedding_theorem}.
\end{proof}

\begin{example}\label{ex:thm:yonedas_embedding_theorem}
  We list examples related to \fullref{thm:yonedas_embedding_theorem}:
  \begin{thmenum}
    \thmitem{ex:thm:yonedas_embedding_theorem/orders} Consider the \hyperref[def:preordered_set]{(pre)ordered set} \( (P, \leq) \), which can be regarded as a \hyperref[def:thin_category]{thin category} in accordance with \fullref{thm:order_category_isomorphism}.

    \Fullref{thm:yonedas_embedding_theorem} implies that, for \( x \) and \( y \) in \( P \),
    \begin{equation*}
      \hom(x, y) \cong \hom(\hom(\anon*, x), \hom(\anon*, y)).
    \end{equation*}

    The set \( \hom(x, y) \) is nonempty if and only if \( x \leq y \).

    The functor \( \hom(\anon*, x) \) can be used to test whether \( \hom(z, x) \) is nonempty for a particular \( z \) in \( P \). A function from \( \hom(z, x) \) to \( \hom(z, y) \) exists unless \( \hom(z, y) \) is empty while \( \hom(z, x) \) is not, i.e. unless \( z \leq x \) but not \( z \leq y \).

    Thus, a natural transformation from \( \hom(\anon*, x) \) to \( \hom(\anon*, y) \) exists if and only if \( z \leq x \) implies \( z \leq y \). Therefore, we can interpret the conclusion of Yoneda's theorem as follows:
    \begin{center}
      We have \( x \leq y \) if and only if, for every \( z \) in \( P \), \( z \leq x \) implies \( z \leq y \).
    \end{center}

    This can also be expressed via initial segments:
    \begin{center}
      We have \( x \leq y \) if and only if \( P_{\leq x} \subseteq P_{\leq y} \).
    \end{center}

    \thmitem{ex:thm:yonedas_embedding_theorem/cayleys_theorem} Consider the \hyperref[def:monoid_delooping]{delooping} \( \cat{B}(M) \) of some \hyperref[def:monoid]{monoid} \( \cat{B}(M) \).

    \Fullref{thm:yonedas_embedding_theorem} implies that
    \begin{equation*}
      \underbrace{\hom(\Anon, \Anon)}_M \cong \hom(\hom(\anon*, \Anon), \hom(\anon*, \Anon)).
    \end{equation*}

    A natural transformation \( \eta: \hom(\anon*, \Anon) \to \hom(\anon*, \Anon) \) consists of a single component
    \begin{equation*}
      \eta_{\Anon}: \hom(\Anon, \Anon) \to \hom(\Anon, \Anon),
    \end{equation*}
    which can be regarded as an endofunction on \( M \). For every monoid element \( x \) (regarded as an endomorphism on \( \Anon \)), the naturality condition implies that
    \begin{equation*}
      \begin{aligned}
        \includegraphics[page=1]{output/ex__thm__yonedas_embedding_theorem}
      \end{aligned}
    \end{equation*}
    that is, for every monoid element \( s \), since \( s \bincirc x = xs \),
    \begin{equation*}
      \eta_{\Anon}(xs) = \eta_{\Anon}(x) s.
    \end{equation*}

    An obvious choice for \( \eta_{\Anon} \) is \( x \mapsto mx \) for some fixed \( m \). This map satisfies the naturality condition due to associativity.

    Yoneda's theorem implies that each of the natural transformations is given by such a map. \Fullref{thm:cayleys_theorem_for_monoids} shows that the monoid structure on the natural transformations is isomorphic to \( M \).

    If \( M \) is a group, every element is invertible and, by \cref{thm:isomorphisms_via_yoneda_embedding}, each natural transformation \( \eta \) is also invertible. Thus, in this case the natural transformations form a group isomorphic to \( G \). This is captured in \fullref{thm:cayleys_theorem}.
  \end{thmenum}
\end{example}

\paragraph{Representable functors}

\begin{definition}\label{def:functor_representation}\mcite[def. 2.1.4]{Riehl2016CategoryTheory}
  Fix a \hyperref[def:small_category]{locally small} \hyperref[def:category]{category} \( \cat{C} \).

  A \term{representation} of the functor \( F: \cat{C} \to \cat{Set} \) is a pair \( (A, \rho) \), where \( A \) is an object in \( \cat{C} \) and \( \rho: \hom(A, \anon*) \to F \) is a \hyperref[def:natural_isomorphism]{natural isomorphism}.

  We can similarly define a representation of a \hyperref[def:presheaf]{presheaf} \( F: \cat{C}^\oppos \to \cat{Set} \) as a pair \( (A, \rho) \), where \( \rho: \hom(\anon*, A) \to F \) is a natural isomorphism.

  In both cases, we call \( F \) \term{representable} if it has a representation.
\end{definition}
\begin{comments}
  \item Here \( \hom(A, \anon*) \) and \( \hom(\anon*, A) \) are the covariant and contravariant \( \hom \)-functors discussed in \cref{def:hom_functor}. If \( \cat{C} \) is not locally small, these functors are ill-defined.

  \item \Fullref{thm:yonedas_lemma} allows characterizing representations --- see \cref{rem:yoneda_representation_characterization}.
\end{comments}

\begin{proposition}\label{thm:def:functor_representation}
  \hyperref[def:functor_representation]{Functor representations} have the following basic properties:
  \begin{thmenum}
    \thmitem{thm:def:functor_representation/uniqueness} All representing objects for the same functor (resp. presheaf) are isomorphic.
  \end{thmenum}
\end{proposition}
\begin{proof}
  \SubProofOf{thm:def:functor_representation/uniqueness} Fix a functor \( F: \cat{C} \to \cat{Set} \) and suppose we have two natural isomorphisms
  \begin{align*}
    &\rho: \hom(A, \anon*) \to F, \\
    &\sigma: \hom(B, \anon*) \to F.
  \end{align*}

  Then, by \cref{thm:def:morphism_invertibility/composition}, the transformation \( \sigma \bincirc \rho^{-1}: \hom(\anon*, A) \to \hom(\anon*, B) \) is also an isomorphism.

  \Cref{thm:isomorphisms_via_yoneda_embedding} implies that \( A \) and \( B \) are isomorphic.

  The case where \( F \) is a presheaf follows by taking \( \cat{C}^\oppos \) instead of \( \cat{C} \).
\end{proof}

\begin{example}\label{ex:def:functor_representation}
  We list examples of \hyperref[def:functor_representation]{functor representations}:
  \begin{thmenum}
    \thmitem{ex:def:functor_representation/global_elements} In \cref{ex:def:natural_transformation/global_elements}, we provide a natural isomorphism
    \begin{equation*}
      \begin{aligned}
        &\alpha: \id_{\cat{Set}} \to \hom(1, \anon*), \\
        &\alpha_A(x) = (0 \mapsto x)
      \end{aligned}
    \end{equation*}
    between set-theoretic elements of a set \( A \) and its \hyperref[def:global_element]{global elements} (morphisms from \( 1 \)).

    Thus, \( (1, \alpha^{-1}) \) is a representation of the identity functor on \( \cat{Set} \).

    \thmitem{ex:def:functor_representation/power_set} As discussed in \cref{ex:def:natural_transformation/boolean}, the \hyperref[def:subset_characteristic_function]{subset characteristic functions} \( \chi_A \) for small sets can be grouped into a natural isomorphism \( \chi: \pow \to \hom(2, \anon*) \).

    Thus, \( (2, \chi^{-1}) \) is a representation of the power set functor on \( \cat{Set} \).
  \end{thmenum}
\end{example}

\begin{remark}\label{rem:yoneda_representation_characterization}\mcite[\S 4.4.2]{Schubert1972Categories}
  In a locally small category \( \cat{C} \), fix a functor \( F: \cat{C} \to \cat{Set} \) and an object \( A \) in \( \cat{C} \). \Fullref{thm:yonedas_lemma} implies that all natural transformations from \( \hom(A, \anon*) \) to \( F(A) \) have the form
  \begin{equation*}
    \alpha_{(A, F)}(x) = \seq[\big]{ \underbrace{s}_{\mathclap{A \to C}} \mapsto \underbrace{F(s)(x)}_{F(C)} }_{C \in \obj(\cat{C})}
  \end{equation*}
  for some \( x \) in \( F(A) \), and are in fact uniquely determined by \( x \).

  This transformation is an isomorphism if and only if, for every object \( C \) in \( \cat{C} \) and every element \( y \) of \( F(C) \), there exists a unique function \( s: A \to C \) such that \( F(s)(x) = y \).

  Thus, the pair \( (A, x) \) determines a \hyperref[def:functor_representation]{representation} \( (A, \rho) \) of \( F \) if and only if this unique existence condition holds.
\end{remark}
